Fix $n\ge 1$.
A set $S\subseteq\{1,2,\dots,n\}$ is \emph{admissible} if
\[
\forall a\ne b\in S,\qquad \gcd(a,b)=1.
\]
Define
\[
G(n)=\max_{\substack{S\subseteq\{1,\dots,n\}\\ S\text{ admissible}}}\ \sum_{a\in S} a.
\]
Define
\[
H(n)=\sum_{\substack{p\text{ prime}\\ p<n}} p\ +\ n\,\pi(\sqrt{n}),
\]
where $\pi(x)=\#\{p\le x: p\text{ prime}\}$.

Questions in the problem text:
\begin{itemize}
\item Is it true that $G(n)>H(n)-n^{1+o(1)}$?
\item Is it true that for every fixed $k\ge 2$, for all sufficiently large $n$, an admissible set achieving $G(n)$ contains an integer with at least $k$ prime factors? (The problem text does not specify whether ``prime factors'' means with multiplicity $\Omega(\cdot)$ or distinct prime factors $\omega(\cdot)$.)
\end{itemize}

The problem text states that Erd\H{o}s and van Lint proved bounds
\[
H(n)-n^{3/2-o(1)}<G(n)<H(n)
\quad\text{and}\quad
\frac{H(n)-G(n)}{n}\to\infty,
\]
and proved the second claim for $k=2$.
